\section{Introdução}

O problema do Corte Máximo (\emph{Max-Cut}) consiste em dividir os vértices de um grafo em dois grupos de modo que o peso total das arestas entre os grupos seja o maior possível. No Trabalho~1 estudamos uma heurística gulosa com busca local 1-flip, que processa grafos de 100 mil vértices em milissegundos, mas não dá nenhuma garantia além de $1/2$ do ótimo e, nos grafos toroidais com pesos $\pm1$ do G-set, ficou em torno de 80\% do melhor valor conhecido. Esse número, porém, é relativo a valores obtidos por metaheurísticas \cite{Benlic and Hao 2013}, e não ao ótimo: não sabíamos o quanto da diferença era culpa da heurística e o quanto vinha de uma referência que também pode estar abaixo do ótimo.

Este trabalho aborda a outra ponta do problema: algoritmos \emph{exatos}, que sempre devolvem o ótimo e provam que ele é o ótimo, ao custo de tempo exponencial no pior caso. Implementamos dois, de técnicas diferentes. O primeiro é um \emph{Branch-and-Bound} com um limitante superior atualizado de forma incremental e inicializado com a heurística do Trabalho~1. O segundo é uma \emph{programação dinâmica sobre a largura de caminho} do grafo, que explora o fato de os grafos esparsos poderem ser processados com uma ``fronteira'' pequena de vértices ativos. Os dois foram validados contra força bruta em 3\,000 grafos aleatórios pequenos.

Em vez de muitas instâncias pequenas, o estudo experimental usa poucas instâncias de tamanho significativo (12, de 100 a 100\,002 vértices) e um limite de 15 minutos por algoritmo em cada uma, o que permite observar até onde cada método chega: em algumas instâncias o ótimo é calculado e comparado com a heurística, em outras o Branch-and-Bound trabalha até o limite de tempo e devolve a melhor solução encontrada. Um resultado negativo precisa ser dito logo: nenhum algoritmo exato escrito do zero resolve as instâncias do G-set usadas no Trabalho~1 (2\,000 a 20\,000 vértices), e a literatura que resolve esses tamanhos recorre a ferramentas muito mais sofisticadas \cite{Liers et al. 2004, Rendl et al. 2010}. Por isso as instâncias combinam três tipos: grafos aleatórios esparsos pequenos, grades toroidais $\pm1$ estreitas (que chegam a 100 mil vértices por terem largura pequena) e instâncias sem estrutura favorável, incluindo duas do G-set (G32 e G57), em que só o Branch-and-Bound roda. A Seção~\ref{sec:instancias} justifica a escolha.

O artigo segue a ordem pedida: as Seções~2 e~3 descrevem o problema, sua complexidade e o estado da arte em algoritmos exatos; a Seção~4 justifica as duas técnicas; a Seção~5 apresenta os algoritmos com provas de correção e análise de complexidade; a Seção~6 descreve as instâncias e a metodologia de teste; a Seção~7 traz os resultados e a comparação com o Trabalho~1; a Seção~8 conclui.

